\documentclass[14pt]{extreport}
\usepackage[T1]{fontenc}
\usepackage[left=1in, right=1in, top=1in, bottom=1in]{geometry}
\usepackage{amsmath, amssymb, amsthm}
\usepackage{enumitem}
\usepackage{mathtools}
\usepackage{cancel}

% colored links
\usepackage{hyperref}
\hypersetup{
    colorlinks=true,
    linkcolor=blue,
    filecolor=magenta,      
    urlcolor=black!20!BurntOrange,
    }

% Inputting Python code
\usepackage[dvipsnames]{xcolor}
\definecolor{textblue}{rgb}{.2,.2,.7}
\definecolor{textred}{rgb}{0.54,0,0}
\definecolor{textgreen}{rgb}{0,0.43,0}
\usepackage{upquote}
\usepackage{listings}
\lstset{
    language=Python, 
    tabsize=4,
    basicstyle={\ttfamily},
    keywordstyle=\color{textblue},
    commentstyle=\color{textgreen},
    stringstyle=\color{textred},
    frame=none,
    columns=fullflexible,
    keepspaces=true,
    showstringspaces=false,
    xleftmargin=-15mm, % manual adjustment, figure out permanent solution
}

% Drawing figures
\usepackage{tikz}
\usetikzlibrary{matrix, positioning}
\usetikzlibrary{calc, shapes.symbols}

% Colored Boxes
\usepackage{tcolorbox}
\tcbuselibrary{skins,hooks}
\usetikzlibrary{shadows}
\usepackage{lipsum}

%Images
\usepackage{graphicx}
    \usepackage{subcaption}
    \usepackage{float}

%Formatting and Spacing
\setitemize[1]{noitemsep, parsep = 5pt, topsep = 5pt}
\setenumerate[1]{label = (\alph*), parsep = 1pt, topsep = 5pt}
\setlength\parindent{0pt}
\linespread{1.1}

%Stats
\newcommand{\Cov}{\text{Cov}}
\newcommand{\trace}{\text{trace}}
\newcommand{\Normal}{\text{Normal}}
\newcommand{\var}[1]{\operatorname{var}\left(#1\right)}
\newcommand{\E}[1]{\operatorname{E}\left[#1\right]}
\renewcommand{\Pr}[1]{\operatorname{P}\left(#1\right)}

%Cases Environment
\newlist{Cases}{enumerate}{3}
\setlist[Cases]{leftmargin = .25in, label = {Case \arabic*.}, topsep = 0.01in, itemsep = 0.04in, itemindent = .5in, parsep = 0in}

%Custom Title Fields
\newcommand{\lectTitle}{Lecture 7 Notes}
\newcommand{\lectTime}{March 7, 2022}
\newcommand{\lectClass}{Advanced Algorithms}
\newcommand{\lectClassInstructor}{Dana Randall}
\newcommand{\lectSection}{Spring 2024}
\newcommand{\lectAuthorName}{Sarthak Mohanty, Aadit Trivedi}

\title{
    \vspace{1.75in}
    \textbf{\lectClass:\\ \lectTitle}\\
    \vspace{0.1in}\large{\textit{\lectClassInstructor\ \lectSection}}
    \vspace{2in}
    \author{\textbf{Sarthak Mohanty} \\ \textbf{Aadit Trivedi}}
    \date{}
}

\begin{document}

\maketitle
\pagebreak

\subsection*{Randomized Min-cut}
    Given some connected graph $G = (V, E)$, a \textbf{cut} partitions the vertices into two sets: $\mathcal{S}$ and $\mathcal{T}$.

    \vspace{4mm}

    One of the motivating questions here is what is the minimum \# of edges whose removal disconnects the graph?

    \vspace{4mm}

    We will not show why here, but one reason why the minimum-cut problem is of interest is that it is tightly connected to the maximum-flow problem. By the max-flow min-cut theorem, the value of the maximum flow from a source $s$ to a sink $t$ is equal to the capacity of the minimum $s$-$t$ cut.

    \vspace{4mm}

    Thus, one natural algorithm for solving the global min-cut problem is as follows: Pick some source vertex $s$, then for every other vertex $t \ne s$, compute the minimum $s$-$t$ cut. Choose the minimum of all these values.

    \vspace{4mm}

    However, this requires $\mathcal{O}(n)$ calls to the max-flow algorithm. Can we do better? In fact, we can! We will now discuss Karger's algorithm, a randomized algorithm for computing the min-cut.

    \vspace{4mm}

    The algorithm is as follows.
    \begin{itemize}
        \item Repeat until two supernodes remain.
        \begin{itemize}
            \item Choose an edge at random $e = (u,v)$.
            \item Contract this edge, making $[uv]$ a supernode. Remove any self loops and keep all multi-edges.
        \end{itemize}
    \end{itemize}

    
    \begin{figure}
        \centering
        \includegraphics{lecture_notes/karger_illustration.png}
        \caption{Karger's Algorithm at work (taken from \href{https://www.cs.princeton.edu/courses/archive/fall13/cos521/lecnotes/lec2final.pdf}{here}.)}
    \end{figure}

    Let us note some observations.
    \begin{itemize}
        \item Note that there can be potentially many min cuts. An example of a graph that has a unique min cut, on the other hand, is two connected graphs connected with a single edge.
        \item The min cut size is at most the minimum degree.
    \end{itemize}

    This algorithm always returns some cut. But now, we propose the question: Is this cut a minimum cut? We argue that this algorithm gives us a minimum cut with some nontrivial probability.

    \vspace{4mm}

    So, what's the probability that we succeed in finding the min cut with this randomized algorithm? Fix some min cut $C$ of size $k$. The only way this fixed cut is destroyed is if we contract one of the $k$ edges crossing $C$.

    \vspace{4mm}

    Since the min cut has size $k$, every vertex has degree at least $k$. Thus, when there are $n$ vertices remaining, there are at least $\frac{kn}{2}$ edges remaining.

    \vspace{4mm}

    What is the probability that $C$ survives the first contraction? This is equal to $1 - \Pr{\text{Failure}}$. The probability of failure is at most
    $$
    \frac{k}{kn/2} = \frac{2}{n}.
    $$
    Thus, the probability that $C$ survives the first contraction is at least
    $$
    1 - \frac{2}{n}.
    $$
    More generally, when there are $i$ supernodes remaining, the contracted graph still has min cut at least $k$, as long as we have not contracted an edge of $C$. Therefore, every supernode has degree at least $k$, so there are at least $\frac{ki}{2}$ edges remaining. Thus, the probability that the next contraction destroys $C$ is at most
    $$
    \frac{k}{ki/2} = \frac{2}{i}.
    $$

    Finally, we can calculate the probability that $C$ survives every contraction, as follows:
    \begin{align*}
        \Pr{\text{$C$ survives}}
        &\geq \prod_{i = 3}^{n} \left(1 - \frac{2}{i}\right) \\
        &= \prod_{i = 3}^{n} \frac{i - 2}{i} \\
        &= \frac{1}{3} \cdot \frac{2}{4} \cdot \frac{3}{5} \cdots \frac{n - 2}{n} \\
        &= \frac{2}{n(n - 1)}.
    \end{align*}

    Thus, one run of Karger's algorithm outputs a fixed minimum cut with probability at least
    $$
    \frac{2}{n(n - 1)} = \frac{1}{\binom{n}{2}}.
    $$

    This also implies that an $n$-vertex graph has at most $\binom{n}{2}$ distinct minimum cuts. Indeed, Karger's algorithm outputs exactly one cut, and every minimum cut is output with probability at least $\frac{1}{\binom{n}{2}}$. Therefore, there can be at most $\binom{n}{2}$ minimum cuts.

\subsection*{Boosting the probability of success}
    Let
    $$
    p_0 \geq \frac{2}{n(n - 1)}
    $$
    be the probability that one run of Karger's algorithm succeeds. If we run the algorithm independently $N$ times and keep the smallest cut found, then the probability of success is
    $$
    P^* = 1 - (1 - p_0)^N.
    $$

    \vspace{4mm}

    Therefore,
    $$
    P^* \geq 1 - \left(1 - \frac{2}{n(n - 1)}\right)^N.
    $$

    \vspace{4mm}

    We use the approximation
    $$
    \lim_{n \to \infty} \left(1 + \frac{\alpha}{f(n)}\right)^{f(n)} = e^\alpha
    $$
    as long as $\lim_{n \to \infty} f(n) = \infty$.

    \vspace{4mm}

    Now we have:
    \begin{align*}
        N &= n(n - 1), & P^* &\approx 1 - e^{-2}, \\
        N &= n(n - 1)\ln(n), & P^* &\approx 1 - e^{-2\ln(n)} = 1 - n^{-2}, \\
        N &= cn(n - 1)\ln(n), & P^* &\approx 1 - e^{-2c\ln(n)} = 1 - n^{-2c}.
    \end{align*}

    Thus, if we repeat the algorithm $cn(n - 1)\ln(n)$ times, then the probability of failure is at most approximately $n^{-2c}$.

    \vspace{4mm}

    \textbf{Runtime:} Using a straightforward implementation, one run takes $O(mn)$ time. Therefore, after
    $$
    N = O(n^2 \log n)
    $$
    independent runs, the total runtime is
    $$
    O(Nmn) = O(mn^3 \log n).
    $$

\subsection*{Karger-Stein computation}

    The Karger-Stein algorithm is an algorithm that computes the minimum cut of a graph with high probability. It improves the original Karger's algorithm by using a recursive approach.

    \vspace{4mm}

    The key idea is that early contractions are relatively safe, while later contractions are much more dangerous. So instead of contracting all the way down to $2$ vertices, we contract only down to roughly $\frac{n}{\sqrt{2}}$ vertices and then recurse twice.

\subsubsection*{Algorithm}
    The algorithm is as follows:

    \begin{enumerate}
        \item If the graph has at most some constant number of vertices, compute the minimum cut directly.
        \item Let
        $$
        t = \left\lceil \frac{|V|}{\sqrt{2}} + 1 \right\rceil.
        $$
        \item Independently run random contraction twice, each time contracting $G$ down to $t$ vertices. Call the resulting graphs $G_1$ and $G_2$.
        \item Recursively apply the algorithm to $G_1$ and $G_2$.
        \item Return the smaller of the two cuts found.
    \end{enumerate}

\subsubsection*{Computation}

    The improvement of the Karger-Stein algorithm comes from its recursive approach. Instead of running one contraction process all the way to $2$ vertices, it branches into two recursive calls after contracting to about $\frac{n}{\sqrt{2}}$ vertices.

    \vspace{4mm}

    The runtime recurrence is
    \begin{equation}
        T(n) = 2T\left(\left\lceil \frac{n}{\sqrt{2}} \right\rceil + 1\right) + O(n^2),
    \end{equation}
    where the $O(n^2)$ term accounts for the random contractions used to reduce the graph size.

    \vspace{4mm}

    To solve this recurrence, note that each recursive call has size roughly $\frac{n}{\sqrt{2}}$, and there are two recursive calls. Thus, each level of recursion costs
    $$
    O(n^2),
    $$
    since
    $$
    2 \left(\frac{n}{\sqrt{2}}\right)^2 = n^2.
    $$
    The depth of recursion is $O(\log n)$, so the total runtime is
    $$
    T(n) = O(n^2 \log n).
    $$

    This is a significant improvement over the simple repeated version of Karger's algorithm, which takes $O(mn^3 \log n)$ time using a straightforward implementation.

    % https://www.cs.toronto.edu/~anikolov/CSC473W20/Lectures/Karger-Stein.pdf

\end{document}